\documentclass[review,12pt,authoryear]{elsarticle}

\usepackage{amsmath,amssymb,amsfonts,amsthm}
\usepackage{booktabs}
\usepackage{graphicx}
\usepackage{hyperref}
\usepackage{setspace}
\usepackage{caption}
\usepackage{subcaption}
\usepackage{microtype}
\usepackage{tabularx}
\usepackage{enumitem}
\usepackage{geometry}
\geometry{margin=1in}

\hypersetup{
    colorlinks=true,
    linkcolor=blue,
    citecolor=blue,
    urlcolor=blue
}

\journal{International Review of Financial Analysis}

\begin{document}

\begin{frontmatter}

\title{Causal Emergence in Financial Markets: Dynamic Organization and Effective Dimensionality During Systemic Stress}

\author{}

\begin{abstract}
This paper investigates whether multi-asset financial systems exhibit causal emergence, such that a macroscopic coarse-graining of the financial network contains higher effective transition information per degree of freedom than its constituent microscopic assets. Applying continuous-state information theory to daily U.S. equity industry portfolio returns from January 1990 to June 2026, we develop the Causal Emergence Financial Index ($\mathrm{CEFI}_t$) and estimate the Causal Effective Dimension ($q_t^*$) via Riemannian optimization on Stiefel manifolds under scale-adaptive Gaussian interventions. Benchmarking against a four-tier hierarchy of matched surrogate null models ($B=9,999$ for primary nulls), we find that causal emergence is state-dependent. In the 2005 calm-market benchmark window, observed emergence is statistically indistinguishable from surrogates preserving own-lag persistence and contemporaneous innovation covariance ($p_{\text{Holm}} = 0.6215$). In contrast, during the 2008 Global Financial Crisis peak and the March 2020 COVID shock, $\mathrm{CEFI}_t$ significantly exceeds both static correlation and cross-lag network isolation nulls ($p_{\text{Holm}} \le 0.0042$), demonstrating structured cross-sectoral dynamical coupling. Systemic liquidity and contagion crises concentrate dynamics into lower macro-dimensions (combined modal $q^* = 2$, with 80.8\% having $q^* \le 4$) relative to valuation repricing episodes (modal $q^* = 4$, with 48.9\% having $q^* \le 4$). Across 870 historical rolling windows, $\mathrm{CEFI}_t$ is positively associated with external continuous-emergence benchmarks ($\rho = 0.837$ and $0.832$). These findings indicate that continuous causal emergence provides an informative measure of collective market organization during systemic stress.
\end{abstract}

\begin{keyword}
Causal Emergence \sep Market Organization \sep Systemic Stress \sep Effective Information \sep Stiefel Manifold \sep Effective Dimensionality
\end{keyword}

\end{frontmatter}

\doublespacing

\section{Introduction}
\label{sec:intro}

The transmission of distress across financial markets has traditionally been studied through the lens of contemporaneous co-movement, cross-sectional factor exposure, and pairwise connectedness \citep{billio2012econometric, diebold2012better, adrian2016covar, brownlees2017srisk}. During periods of severe market instability, asset returns frequently exhibit heightened cross-sectional correlation and variance concentration, phenomena commonly described as financial contagion or systemic fragility \citep{acharya2017measuring, kritzman2011principal}. However, standard econometrics evaluates financial dynamics predominantly at the observed microscopic asset resolution or through unconditional factor projections. It does not address a distinct architectural question: \emph{does a distressed financial market develop macroscopic dynamical organization, such that a coarse-grained representation of the network possesses greater effective transition information per degree of freedom than its underlying constituent assets?}

The financial significance of this question lies in distinguishing between two fundamentally different market stress mechanisms: (i) common simultaneous market stress, where asset returns experience contemporaneous shocks driven by common factors without structured lead-lag transmission; and (ii) organized lagged propagation across sectors, where distress cascades through intertemporal feedback channels. While static covariance metrics and absorption ratios capture contemporaneous variance concentration, they cannot determine whether intertemporal transitions concentrate on a low-dimensional manifold that enhances dynamical predictability relative to noisy microscopic returns.

In the physics of complex systems and information theory, the formal framework of \emph{Causal Emergence} (CE), introduced by \citet{hoel2013quantifying} and extended by \citet{klein2020emergence} and \citet{rosas2020reconciling}, demonstrates that macroscopic representations of noisy or degenerate networks can exhibit higher \emph{Effective Information} ($EI$) than their microscopic descriptions. By grouping or projecting micro-level states onto an optimal lower-dimensional subspace, macro coarse-graining can average out idiosyncratic noise while concentrating intertemporal transition dynamics. Recent mathematical advances have established exact analytical formulations of causal emergence for continuous-state linear stochastic iterative systems \citep{liu2024exact, liu2025singular, yang2025finding}.

This paper develops a rolling empirical implementation of continuous-state causal emergence using daily U.S. industry portfolio returns from January 1990 to June 2026. We model the cross-section of industry portfolio returns through continuous-state linear Gaussian Vector Autoregressive transition models (VAR(1)). Using Riemannian gradient ascent on Stiefel manifolds, we optimize an orthogonal projection matrix $\mathbf{W} \in \mathcal{V}_q(\mathbb{R}^p)$ that maximizes macro effective information across all potential dimensions $q \in \{1, \dots, p-1\}$. We define two central statistics:
\begin{enumerate}[label=(\roman*)]
    \item The \emph{Causal Emergence Financial Index} ($\mathrm{CEFI}_t$), measuring the maximal causal information rate per degree of freedom gained through macroscopic coarse-graining:
    \begin{equation}
    \mathrm{CEFI}_t = \max_{1 \le q < p} \left[ \frac{EI_{q,t}^*}{q} - \frac{EI_{p,t}}{p} \right]
    \end{equation}
    \item The \emph{Causal Effective Dimension} ($q_t^*$), identifying the macro dimension at which causal concentration is maximized:
    \begin{equation}
    q_t^* = \arg\max_{1 \le q < p} \left[ \frac{EI_{q,t}^*}{q} - \frac{EI_{p,t}}{p} \right]
    \end{equation}
\end{enumerate}

To ensure that the measure is invariant to common scalar changes in return units (e.g., decimal vs. percentage returns), we introduce an isotropic Gaussian intervention scaled by average channel energy, combined with an energy-scaled ridge regularizer. Because positive $\mathrm{CEFI}_t$ can arise mechanically under finite-sample optimization, we evaluate historical estimates against a four-tier hierarchy of matched surrogate null models ($B=9,999$ for primary nulls and $B=999$ for auxiliary nulls), including a cross-lag network isolation null ($H_0^{\text{diag+contemp}}$) that preserves own-lag persistence and the full empirical contemporaneous innovation covariance matrix $\boldsymbol{\Sigma}_\varepsilon$.

The empirical analysis delivers several main findings:
\begin{enumerate}
    \item \textbf{State- and Episode-Dependent Emergence:} In the 2005 calm-market benchmark window, observed $\mathrm{CEFI}_t$ is statistically indistinguishable from surrogate benchmarks preserving own-lag persistence and contemporaneous covariance ($p_{\text{emp}} = 0.6215, p_{\text{Holm}} = 0.6215$). In contrast, during the 2008 Global Financial Crisis (GFC) peak and the March 2020 COVID shock, observed $\mathrm{CEFI}_t$ significantly exceeds the static correlation mode benchmark ($p_{\text{emp}} = 0.0014, p_{\text{Holm}} = 0.0042$ and $p_{\text{emp}} = 0.0001, p_{\text{Holm}} = 0.0006$, respectively).
    \item \textbf{Cross-Lag Dynamical Organization During Systemic Stress:} At the peak of the 2008 GFC and the 2020 COVID shock trough, observed $\mathrm{CEFI}_t$ significantly exceeds the cross-lag network isolation null ($p_{\text{emp}} = 0.0001, p_{\text{Holm}} = 0.0006$), demonstrating structured intertemporal organization that cannot be explained by static covariance and own-lag persistence alone.
    \item \textbf{Dimensionality Concentration:} Systemic liquidity and contagion episodes exhibit greater concentration toward low-dimensional macro representations, with a combined modal $q^*=2$ and 80.8\% of observations satisfying $q^* \le 4$, compared with 48.9\% during valuation repricing episodes (modal $q^*=4$).
    \item \textbf{Cross-Method Benchmarking:} Across 870 historical rolling windows, $\mathrm{CEFI}_t$ is positively associated over time with the \citet{liu2024exact} and \citet{liu2025singular} benchmarks (Pearson $\rho = 0.837$ and $0.832$, respectively). Dimensional agreement is high for the SVD benchmark (88.7\% within $\pm 1$) but low for the uniform-intervention formulation.
    \item \textbf{Cross-Universe Robustness:} On the 49 Fama-French industry cross-section ($p=49$), the pipeline yields a modal dimension of $q^* = 3$, with 71.7\% of historical windows exhibiting $q^* \le 4$.
\end{enumerate}

The remainder of the paper is organized as follows. Section \ref{sec:literature} reviews the related literature. Section \ref{sec:methodology} presents the mathematical methodology and Stiefel optimization. Section \ref{sec:data} describes the data and matched null hierarchy. Section \ref{sec:results} reports empirical findings. Section \ref{sec:robustness} provides theoretical benchmarking and robustness checks. Section \ref{sec:discussion} discusses economic interpretations. Section \ref{sec:limitations} outlines limitations, and Section \ref{sec:conclusion} concludes.

\section{Related Literature and Conceptual Foundations}
\label{sec:literature}

\subsection{Systemic Risk and Financial Connectedness}
Our work relates to the econometric literature on systemic financial risk, market connectedness, and cross-sectoral spillover networks \citep{billio2012econometric, diebold2012better, diebold2014network, adrian2016covar, brownlees2017srisk}. Recent contributions in the \emph{International Review of Financial Analysis} and related finance literature have highlighted the importance of dynamic connectedness, quantile risk transmission, and sector fragility during crisis episodes \citep{ahelegbey2022network, yousaf2022linkages, mensi2021spillovers, bouri2021return, yang2023modeling}. 

While these frameworks effectively quantify contemporaneous co-movement, tail dependence, or pairwise linear spillovers at the observed microscopic resolution, they do not examine whether a macroscopic aggregation of the system can exhibit higher deterministic causal transition structure than the underlying portfolio components.

\subsection{Information Theory, Entropy, and Dimensionality in Finance}
Information-theoretic concepts have a rich history in quantitative finance. Non-parametric measures such as Shannon entropy, transfer entropy, and mutual information have been applied to detect non-linear information flows and cross-market spillover networks \citep{schreiber2000measuring, dimpfl2013using, bekiros2017information}. Spectral measures of effective dimensionality, such as the effective rank of a covariance matrix \citep{roy2007effective}, have been employed to quantify market concentration during crises.

However, existing spectral dimensionality metrics depend exclusively on the static contemporaneous covariance matrix $\boldsymbol{\Sigma}_x$. They do not incorporate intertemporal transition dynamics ($\mathbf{A}$) or innovation noise covariance ($\boldsymbol{\Sigma}_\varepsilon$). In contrast, the Causal Effective Dimension ($q_t^*$) developed in this paper incorporates transition mechanics and innovation uncertainty simultaneously under an interventional perturbation framework.

\subsection{Causal Emergence in Complex Dynamical Systems}
The theory of causal emergence originated in discrete Markov chains to formalize how coarse-grained macroscopic states can reduce degeneracy and noise \citep{hoel2013quantifying, hoel2017when, klein2020emergence}. \citet{rosas2020reconciling} unified causal emergence with partial information decomposition, distinguishing between downward causation and synergistic emergence.

Recent mathematical physics research has extended causal emergence to continuous-state stochastic dynamical systems. \citet{liu2024exact} derived exact continuous effective information and emergence conditions for linear stochastic systems under uniform bounded interventions. \citet{liu2025singular} developed an analytical singular value decomposition (SVD) perspective for Gaussian iterative dynamics, demonstrating equivalence conditions between SVD-based and EI-based emergence. \citet{yang2025finding} formulated machine learning architectures to discover emergence from observational data. Our paper adapts these continuous-state foundations to empirical financial econometrics, rolling-window market monitoring, and matched-null inference.

\section{Mathematical Methodology}
\label{sec:methodology}

\subsection{Reduced-Form Linear Continuous-State Markov Dynamics}
Let $\mathbf{x}_t = (x_{1,t}, x_{2,t}, \dots, x_{p,t})^\top \in \mathbb{R}^p$ denote a $p$-dimensional column vector of daily portfolio returns. We model the transition dynamics via a continuous-state first-order Vector Autoregression, $\text{VAR}(1)$:
\begin{equation}
\mathbf{x}_{t+1} = \mathbf{A}\mathbf{x}_t + \boldsymbol{\varepsilon}_{t+1}, \qquad \boldsymbol{\varepsilon}_{t+1} \sim \mathcal{N}(\mathbf{0}, \boldsymbol{\Sigma}_\varepsilon)
\label{eq:var1_model}
\end{equation}
where $\mathbf{A} \in \mathbb{R}^{p \times p}$ is the transition matrix, and $\boldsymbol{\varepsilon}_{t+1}$ is a vector of Gaussian innovations with positive-definite covariance matrix $\boldsymbol{\Sigma}_\varepsilon \in \mathbb{R}^{p \times p}$. 

Let $\mathbf{X}_{\text{lag}} \in \mathbb{R}^{(T-1) \times p}$ and $\mathbf{X}_{\text{lead}} \in \mathbb{R}^{(T-1) \times p}$ denote the data matrices whose $t$-th rows contain $\mathbf{x}_t^\top$ and $\mathbf{x}_{t+1}^\top$, respectively. In row-regression format, $\mathbf{X}_{\text{lead}} \approx \mathbf{X}_{\text{lag}}\mathbf{A}^\top$. To guarantee scale invariance and prevent overfitting in finite rolling estimation windows, $\mathbf{A}$ is estimated via trace-scaled ridge regression:
\begin{equation}
\hat{\mathbf{A}} = \mathbf{X}_{\text{lead}}^\top \mathbf{X}_{\text{lag}} \left( \mathbf{X}_{\text{lag}}^\top \mathbf{X}_{\text{lag}} + \lambda_t \mathbf{I}_p \right)^{-1}, \qquad \lambda_t = \lambda_0 \cdot \frac{\operatorname{Tr}(\mathbf{X}_{\text{lag}}^\top \mathbf{X}_{\text{lag}})}{p}
\label{eq:scaled_ridge}
\end{equation}
with baseline regularization parameter $\lambda_0 = 10^{-4}$. Scaling the ridge penalty proportionally to the empirical trace of the Gram matrix ensures that $\hat{\mathbf{A}}$ is strictly invariant under common scalar rescalings of return units ($\mathbf{X} \to c\mathbf{X}$). The residual covariance matrix $\boldsymbol{\Sigma}_\varepsilon$ is estimated using Ledoit-Wolf analytical shrinkage \citep{ledoit2004well}, which preserves scale equivariance ($\boldsymbol{\Sigma}_\varepsilon \to c^2 \boldsymbol{\Sigma}_\varepsilon$) and guarantees strict positive definiteness ($\lambda_{\min}(\boldsymbol{\Sigma}_\varepsilon) > 0$). The unconditional covariance matrix of the state vector is denoted by $\boldsymbol{\Sigma}_x = \operatorname{Cov}(\mathbf{x}_t)$.

\paragraph{Delimitation of Causal Terminology.}
The term \emph{causal} is used throughout this paper strictly in the intervention-based information-theoretic sense of Effective Information conditional on the estimated stochastic transition operator \citep{hoel2013quantifying, liu2024exact}. It should not be interpreted as structural economic causality in the sense of Structural VAR (SVAR) identification, instrumental variables, potential outcomes, or Pearlian structural causal DAGs. The estimated VAR transition matrix $\mathbf{A}$ provides a reduced-form dynamical representation of market co-movements upon which counterfactual isotropic Gaussian intervention distributions are evaluated.

\subsection{Isotropic Scale-Invariant Gaussian Interventions and Effective Information}
Effective Information ($EI$) quantifies the mutual information between an intervened state $do(\mathbf{x}_t)$ and the subsequent state $\mathbf{x}_{t+1}$ \citep{hoel2013quantifying}:
\begin{equation}
EI(\mathbf{x}) = I(do(\mathbf{x}_t); \mathbf{x}_{t+1}) = H(\mathbf{x}_{t+1} \mid do) - H(\mathbf{x}_{t+1} \mid \mathbf{x}_t)
\label{eq:ei_def_eq}
\end{equation}
Under Gaussian transition dynamics, the conditional entropy is given by $H(\mathbf{x}_{t+1} \mid \mathbf{x}_t) = \frac{p}{2}\ln(2\pi e) + \frac{1}{2}\ln\det(\boldsymbol{\Sigma}_\varepsilon)$.

To evaluate interventional causality independently of observational factor orientations, we apply an isotropic Gaussian intervention $do(\mathbf{x}_t) \sim \mathcal{N}(\mathbf{0}, \sigma_{do,t}^2 \mathbf{I}_p)$. Conditional on a fixed isotropic covariance, the Gaussian distribution maximizes differential entropy. To ensure that $EI$ is strictly invariant to common scalar changes in return units, we set the intervention variance proportionally to the average channel energy:
\begin{equation}
\sigma_{do,t}^2 = \kappa^2 \cdot \frac{\operatorname{Tr}(\boldsymbol{\Sigma}_{x,t})}{p}
\label{eq:scale_inv_def}
\end{equation}
where $\kappa > 0$ is a dimensionless scale parameter ($\kappa = 1.0$ by default). Under this intervention, $\mathbf{x}_{t+1} \mid do \sim \mathcal{N}(\mathbf{0}, \sigma_{do,t}^2 \mathbf{A}\mathbf{A}^\top + \boldsymbol{\Sigma}_\varepsilon)$. Applying Sylvester's determinant theorem yields the closed-form continuous Effective Information in symmetric positive definite form:
\begin{equation}
EI_{\text{micro}}(\mathbf{x}) = \frac{1}{2} \ln \det\left( \mathbf{I}_p + \sigma_{do,t}^2 \boldsymbol{\Sigma}_\varepsilon^{-1/2} \mathbf{A}\mathbf{A}^\top \boldsymbol{\Sigma}_\varepsilon^{-1/2} \right) = \frac{1}{2} \ln \det\left( \mathbf{I}_p + \sigma_{do,t}^2 \mathbf{A}\mathbf{A}^\top \boldsymbol{\Sigma}_\varepsilon^{-1} \right)
\label{eq:ei_micro_closed}
\end{equation}

\subsection{Macro Coarse-Graining and Stiefel Manifold Optimization}
We define a linear coarse-graining mapping from the $p$-dimensional micro-system to a $q$-dimensional macro-system ($q < p$) via an orthogonal projection matrix $\mathbf{W} \in \mathbb{R}^{q \times p}$:
\begin{equation}
\mathbf{y}_t = \mathbf{W}\mathbf{x}_t, \qquad \mathbf{W}\mathbf{W}^\top = \mathbf{I}_q
\end{equation}
The matrix $\mathbf{W}$ resides on the row-Stiefel manifold $\mathcal{V}_q(\mathbb{R}^p) = \{ \mathbf{W} \in \mathbb{R}^{q \times p} : \mathbf{W}\mathbf{W}^\top = \mathbf{I}_q \}$. 

\paragraph{Interventional Macro Channel Construction.}
Following \citet{liu2024exact}, an intervention on the macro-state $do(\mathbf{y}_t)$ is lifted to the micro-state via the canonical right inverse $do(\mathbf{x}_t) = \mathbf{W}^\top do(\mathbf{y}_t)$. Projecting the resulting transition back via $\mathbf{W}$ defines the constructed macro interventional channel:
\begin{equation}
\mathbf{y}_{t+1} = \mathbf{A}_M \mathbf{y}_t + \boldsymbol{\varepsilon}_{M, t+1}, \qquad \mathbf{A}_M = \mathbf{W}\mathbf{A}\mathbf{W}^\top, \quad \boldsymbol{\Sigma}_M = \mathbf{W}\boldsymbol{\Sigma}_\varepsilon \mathbf{W}^\top
\label{eq:macro_dynamics_def}
\end{equation}
The macro Effective Information is:
\begin{equation}
EI_q(\mathbf{W}) = \frac{1}{2} \ln \det\left( \mathbf{I}_q + \sigma_{do,t}^2 (\mathbf{W}\mathbf{A}\mathbf{W}^\top)(\mathbf{W}\mathbf{A}\mathbf{W}^\top)^\top (\mathbf{W}\boldsymbol{\Sigma}_\varepsilon \mathbf{W}^\top)^{-1} \right)
\label{eq:ei_macro_closed}
\end{equation}

The optimal coarse-graining $\mathbf{W}_q^* = \arg\max_{\mathbf{W} \in \mathcal{V}_q(\mathbb{R}^p)} EI_q(\mathbf{W})$ is sought via Riemannian gradient ascent \citep{edelman1998geometry, absil2008optimization}. Under the canonical metric on $T_{\mathbf{W}}\mathcal{V}_q(\mathbb{R}^p)$, the Riemannian gradient is:
\begin{equation}
\operatorname{grad}_{\mathcal{R}} EI_q(\mathbf{W}) = \mathbf{G} - \mathbf{W}\mathbf{G}^\top \mathbf{W}
\end{equation}
where $\mathbf{G} = \nabla_{\mathbf{W}} EI_q(\mathbf{W})$ is the Euclidean gradient. Updates are mapped back to the Stiefel manifold using QR retraction:
\begin{equation}
\mathbf{W}^{(k+1)} = \operatorname{Retr}\left( \mathbf{W}^{(k)} + \eta \operatorname{grad}_{\mathcal{R}} EI_q(\mathbf{W}^{(k)}) \right) = \mathbf{Q}^\top, \quad \mathbf{Q}\mathbf{R} = \left(\mathbf{W}^{(k)} + \eta \operatorname{grad}_{\mathcal{R}} EI\right)^\top
\end{equation}
where $\eta = 0.05$. We execute 35 gradient iterations across four deterministic multistart initializations (including PCA projection and deterministic orthogonal bases) and retain the solution with the highest objective value attained across the multistart solutions.

\subsection{Causal Emergence Financial Index, Effective Dimension, and Tie-Breaking}
Because continuous $EI$ scales with state dimension, comparisons across dimensions evaluate the \emph{information rate per degree of freedom} \citep{liu2024exact}:
\begin{equation}
\mathrm{CEFI}_t = \max_{1 \le q < p} \left[ \frac{EI_{q,t}^*}{q} - \frac{EI_{p,t}}{p} \right]
\label{eq:cefi_def}
\end{equation}
\begin{equation}
q_t^* = \arg\max_{1 \le q < p} \left[ \frac{EI_{q,t}^*}{q} - \frac{EI_{p,t}}{p} \right]
\label{eq:q_star_def}
\end{equation}
When multiple candidate dimensions $q$ achieve objective values within numerical tolerance ($\epsilon = 10^{-7}$), a deterministic tie-breaking rule selects the minimal dimension $q^* = \min \{ q : \Delta EI_q \ge \max_{k} \Delta EI_k - \epsilon \}$.

\section{Data and Empirical Design}
\label{sec:data}

\subsection{Equity Portfolio Datasets}
Our empirical analysis uses daily value-weighted portfolio returns from the Kenneth French Data Library. The primary dataset consists of the 30 Industry Portfolios (FF30) from January 1990 to June 2026 ($T = 9,190$ trading days, $p = 30$). To evaluate cross-universe robustness, we replicate the pipeline on the 49 Industry Portfolios (FF49, $p=49$).

\subsection{Rolling Estimation Protocol}
We employ rolling estimation windows of $W = 500$ trading days ($\sim 2$ years), stepping by 2 trading days across the 1992--2026 period ($N = 4,346$ rolling parameter estimates). At each step $t$, the trailing window $X_{t-W+1:t}$ is used to estimate $\hat{\mathbf{A}}_t$, $\boldsymbol{\Sigma}_{\varepsilon,t}$, and $\boldsymbol{\Sigma}_{x,t}$, and solve the Stiefel optimization problem for all $q \in \{1, \dots, p-1\}$.

\subsection{Benchmark Historical Episodes and Trailing Window Endpoints}
We identify four major market stress episodes:
\begin{enumerate}
    \item \textbf{2008 Global Financial Crisis (GFC)}: October 1, 2007 to June 30, 2009 (systemic banking contagion and credit freeze).
    \item \textbf{2020 COVID Shock}: February 1, 2020 to May 31, 2020 (simultaneous global market crash and liquidity squeeze).
    \item \textbf{2000 Dot-Com Crash}: March 1, 2000 to October 31, 2002 (valuation bubble collapse concentrated in technology and telecom).
    \item \textbf{2022 Rate Tightening}: January 1, 2022 to November 30, 2022 (macro-monetary discount rate repricing).
\end{enumerate}
Episodes (1) and (2) represent systemic liquidity and contagion crises, whereas (3) and (4) represent broad valuation and rate repricing shocks.

For matched-null inference, benchmark statistics are computed from 500-trading-day trailing windows ending on December 30, 2005, November 20, 2008, and March 23, 2020, corresponding respectively to the calm benchmark (calendar year-end 2005), the peak post-Lehman market liquidation week, and the global market trough of the COVID-19 pandemic selloff. Benchmark endpoints were defined prior to surrogate estimation using externally dated market stress milestones rather than post-hoc optimization of $\mathrm{CEFI}_t$.

\subsection{Hierarchy of Four Matched Surrogate Null Models}
To assess whether empirical $\mathrm{CEFI}_t$ reflects dynamic structure rather than finite-sample optimization artifacts, we evaluate four surrogate null models:
\begin{enumerate}
    \item \textbf{$H_0^{\text{circ}}$ (Circular Shift Null)}: Circularly shifts each asset column independently by random lags. Preserves marginal autocorrelation and variance, destroys cross-correlations and cross-lags.
    \item \textbf{$H_0^{\text{diag}}$ (Decoupled Diagonal VAR Null)}: Simulates an uncoupled VAR(1) with $\mathbf{D} = \operatorname{diag}(\hat{\mathbf{A}})$ and $\boldsymbol{\Sigma}_{\text{diag}} = \operatorname{diag}(\boldsymbol{\Sigma}_\varepsilon)$.
    \item \textbf{$H_0^{\text{static}}$ (Static Correlation Mode Null)}: Synchronously permutes time steps across all assets. Preserves the exact contemporaneous covariance $\boldsymbol{\Sigma}_x$, but destroys intertemporal causality $\mathbf{x}_t \to \mathbf{x}_{t+1}$.
    \item \textbf{$H_0^{\text{diag+contemp}}$ (Cross-Lag Network Isolation Null)}: Simulates $\mathbf{x}_{t+1} = \operatorname{diag}(\hat{\mathbf{A}})\mathbf{x}_t + \boldsymbol{\varepsilon}_{t+1}$ with the full empirical contemporaneous covariance $\boldsymbol{\Sigma}_\varepsilon$. Preserves univariate persistence and contemporaneous shock covariance, destroying only off-diagonal cross-lag couplings ($A_{ij} = 0, i \neq j$).
\end{enumerate}
For each regime, the observed statistic $\mathrm{CEFI}_{\text{obs}}$ is computed once with deterministic initialization and compared against $B=9,999$ surrogate replications for primary nulls ($H_0^{\text{static}}$ and $H_0^{\text{diag+contemp}}$). Empirical $p$-values are computed as $p_{\text{emp}} = (1 + \sum_{b=1}^B \mathbb{I}(\mathrm{CEFI}_b^{\text{null}} \ge \mathrm{CEFI}_{\text{obs}}))/(B + 1)$.

\section{Empirical Results}
\label{sec:results}

\subsection{Evidence from Matched Null Models and Mechanism Contrast}
Table \ref{tab:null_models} reports the statistical inference across calm and crisis episodes under $B=9,999$ surrogate replications for primary nulls, along with Holm-Bonferroni family-wise adjusted $p$-values ($m=6$).

\begin{table}[htbp]
\centering
\caption{Inference on Causal Emergence Across 4 Matched Null Models ($B=9,999$ for Primary Nulls, $B=999$ for Auxiliary Nulls, $q \in 1..29$)}
\label{tab:null_models}
\resizebox{\textwidth}{!}{
\begin{tabular}{llcccccccc}
\toprule
\textbf{Market Regime} & \textbf{Null Model} & $\mathbf{B}$ & $\mathbf{CEFI_{\text{obs}}}$ & $\mathbb{E}[\mathbf{CEFI_0}]$ & $\mathbf{Q_{95}}$ & $\mathbf{z_{\text{dev}}}$ & $\mathbf{p_{\text{emp}}}$ & $\mathbf{p_{\text{Holm}}}$ & \textbf{Modal } $\mathbf{q_0^*}$ \\
\midrule
\textbf{Calm Period (2005)} 
 & $H_0^{\text{circ}}$ (Circular) & 999 & 0.9848 & +0.0280 & +0.0616 & +45.98 & 0.0010 & -- & 29 \\
 & $H_0^{\text{diag}}$ (VAR Diag) & 999 & 0.9848 & +0.0123 & +0.0171 & +350.34 & 0.0010 & -- & 29 \\
 & $H_0^{\text{static}}$ (Static Mode) & 9,999 & 0.9848 & +0.8587 & +1.1494 & +0.85 & 0.1940 & 0.3880 & 2 \\
 & $H_0^{\text{diag+contemp}}$ (Cross-lag Isol.) & 9,999 & 0.9848 & +1.0255 & +1.2228 & -0.35 & 0.6215 & 0.6215 & 3 \\
\midrule
\textbf{2008 GFC Peak} 
 & $H_0^{\text{circ}}$ (Circular) & 999 & 1.8948 & +0.0618 & +0.1622 & +35.35 & 0.0010 & -- & 29 \\
 & $H_0^{\text{diag}}$ (VAR Diag) & 999 & 1.8948 & +0.0583 & +0.1152 & +53.92 & 0.0010 & -- & 1 \\
 & $H_0^{\text{static}}$ (Static Mode) & 9,999 & 1.8948 & +1.3712 & +1.6857 & +3.16 & 0.0014 & 0.0042 & 2 \\
 & $H_0^{\text{diag+contemp}}$ (Cross-lag Isol.) & 9,999 & 1.8948 & +1.4832 & +1.6437 & +4.57 & 0.0001 & 0.0006 & 3 \\
\midrule
\textbf{2020 COVID Shock} 
 & $H_0^{\text{circ}}$ (Circular) & 999 & 1.8787 & +0.0334 & +0.1058 & +50.47 & 0.0010 & -- & 1 \\
 & $H_0^{\text{diag}}$ (VAR Diag) & 999 & 1.8787 & +0.0015 & +0.0035 & +966.78 & 0.0010 & -- & 29 \\
 & $H_0^{\text{static}}$ (Static Mode) & 9,999 & 1.8787 & +1.3154 & +1.5831 & +2.95 & 0.0001 & 0.0006 & 2 \\
 & $H_0^{\text{diag+contemp}}$ (Cross-lag Isol.) & 9,999 & 1.8787 & +1.3977 & +1.5197 & +6.22 & 0.0001 & 0.0006 & 3 \\
\bottomrule
\end{tabular}
}
\end{table}

\begin{figure}[htbp]
\centering
\includegraphics[width=0.95\textwidth]{reports/figures/figure1_cefi_dynamics.pdf}
\caption{Historical dynamics of the Causal Emergence Financial Index ($\mathrm{CEFI}_t$, 1992--2026). Shaded areas indicate historical benchmark stress episodes.}
\label{fig:cefi_dynamics}
\end{figure}

The results in Table \ref{tab:null_models} and Figure \ref{fig:cefi_dynamics} establish that:
\begin{enumerate}
    \item In the 2005 calm-market benchmark window, observed $\mathrm{CEFI}_t = 0.9848$ ($q^* = 3$) is consistent with univariate persistence and contemporaneous shock covariance ($p_{\text{emp}} = 0.6215, p_{\text{Holm}} = 0.6215$).
    \item During the 2008 GFC peak (November 2008), observed $\mathrm{CEFI}_t = 1.8948$ ($q^* = 2$) significantly exceeds both the static correlation mode null ($p_{\text{emp}} = 0.0014, p_{\text{Holm}} = 0.0042$) and the cross-lag network isolation null ($p_{\text{emp}} = 0.0001, p_{\text{Holm}} = 0.0006$), indicating structured off-diagonal dynamical organization across sectors.
    \item During the March 2020 COVID shock, observed $\mathrm{CEFI}_t = 1.8787$ ($q^* = 2$) similarly exceeds both static correlation and cross-lag isolation ($p_{\text{emp}} = 0.0001, p_{\text{Holm}} = 0.0006$).
\end{enumerate}

\begin{figure}[htbp]
\centering
\includegraphics[width=0.95\textwidth]{reports/figures/figure3_null_distributions.pdf}
\caption{Observed $\mathrm{CEFI}_{\text{obs}}$ (red vertical line) compared against moment-matched surrogate distributions of $H_0^{\text{static}}$ and $H_0^{\text{diag+contemp}}$ across benchmark regimes ($B=9,999$). Empirical $p$-values reported in Table \ref{tab:null_models} are exact non-parametric proportions from the full surrogate simulation.}
\label{fig:null_dist}
\end{figure}

\subsection{Response Across Historical Episodes: Descriptive Regime Evidence}
To examine whether $\mathrm{CEFI}_t$ behaves differently across historical market stress episodes, we report disaggregated episode-level summary statistics in Table \ref{tab:episode_summary}.

\begin{table}[htbp]
\centering
\caption{Disaggregated Historical Episode Summary Statistics}
\label{tab:episode_summary}
\resizebox{\textwidth}{!}{
\begin{tabular}{llcccccccc}
\toprule
\textbf{Episode} & \textbf{Stress Class} & \textbf{Obs} & \textbf{Mean } $\mathbf{CEFI}$ & \textbf{Median } $\mathbf{CEFI}$ & \textbf{Mean } $\mathbf{q^*}$ & \textbf{Modal } $\mathbf{q^*}$ & $\mathbf{P(q^* \le 4)}$ & \textbf{Mean } $\mathbf{\bar{\rho}}$ & \textbf{Mean } $\mathbf{ER}$ \\
\midrule
Dot-Com Crash (2000--02)  & Valuation Repricing & 335 & 0.4805 & 0.4391 & 4.93 & 4 & 49.9\% & 0.3387 & 24.60 \\
2008 GFC (2007--09)       & Liquidity/Contagion & 220 & 1.2590 & 1.2272 & 3.55 & 4 & 77.7\% & 0.6833 & 19.36 \\
2020 COVID (2020)         & Liquidity/Contagion & 41  & 1.4706 & 1.6380 & 2.41 & \textbf{2} & \textbf{97.6\%} & 0.6611 & 20.26 \\
2022 Tightening (2022)    & Valuation Repricing & 115 & 0.9546 & 0.8680 & 4.67 & 2 & 46.1\% & 0.5677 & 20.25 \\
\midrule
Calm / Other Periods      & Baseline Benchmark  & 3,635 & 0.9423 & 0.8845 & 3.45 & 3 & 79.2\% & 0.4120 & 23.15 \\
\bottomrule
\end{tabular}
}
\end{table}

We estimate the descriptive event-study regression specification:
\begin{equation}
\mathrm{CEFI}_t = \alpha + \beta_{\text{Liq}} D_t^{\text{Liquidity}} + \beta_{\text{Val}} D_t^{\text{Valuation}} + u_t
\label{eq:h2_reg_spec}
\end{equation}
where $D_t^{\text{Liquidity}} = 1$ during the 2008 GFC and March 2020 COVID crash, and $D_t^{\text{Valuation}} = 1$ during the 2000 Dot-Com crash and 2022 rate tightening.

\begin{table}[htbp]
\centering
\caption{Descriptive Historical Regime Regressions: Sensitivity Across Extended HAC Lag Bandwidths ($L = 20$ to $L = 250$)}
\label{tab:h2_hac}
\resizebox{\textwidth}{!}{
\begin{tabular}{cccccccc}
\toprule
\textbf{HAC Lag ($L$)} & $\mathbf{\beta_{\text{Liq}}}$ & $\mathbf{t\text{-stat}}$ & $\mathbf{\beta_{\text{Val}}}$ & $\mathbf{t\text{-stat}}$ & $\mathbf{\Delta \beta}$ & \textbf{Wald } $\mathbf{t\text{-stat}}$ & \textbf{Wald } $\mathbf{p\text{-val}}$ \\
\midrule
$L = 20$  & +0.336 & +5.50 & -0.355 & -6.30 & +0.691 & +8.31 & $1.11 \times 10^{-16}$ \\
$L = 40$  & +0.336 & +4.54 & -0.355 & -4.82 & +0.691 & +7.11 & $1.16 \times 10^{-12}$ \\
$L = 60$  & +0.336 & +4.07 & -0.355 & -4.15 & +0.691 & +6.24 & $4.38 \times 10^{-10}$ \\
$L = 120$ & +0.336 & +3.64 & -0.355 & -3.28 & +0.691 & +5.20 & $1.99 \times 10^{-07}$ \\
$L = 250$ & +0.336 & +4.04 & -0.355 & -2.89 & +0.691 & +4.98 & $6.36 \times 10^{-07}$ \\
\bottomrule
\end{tabular}
}
\end{table}

Table \ref{tab:h2_hac} shows that across the four benchmark episodes analyzed, $\mathrm{CEFI}_t$ was higher during liquidity and contagion dislocations ($\beta_{\text{Liq}} = +0.336$) than during valuation repricing episodes ($\beta_{\text{Val}} = -0.355$), yielding an estimated contrast of $\Delta \beta = +0.691$. In leave-one-episode-out sensitivity checks (Supplementary Table A12), $\Delta \beta$ remains strictly positive ($\Delta \beta \in [+0.338, +0.869]$, Wald $t \ge 3.10$) upon removing each episode sequentially.

\subsection{Behavior of Causal Effective Dimensionality Across Market Regimes}
Table \ref{tab:episode_summary} and Figure \ref{fig:qstar_dynamics} report the distribution of $q_t^*$.

\begin{figure}[htbp]
\centering
\includegraphics[width=0.95\textwidth]{reports/figures/figure2_qstar_dynamics.pdf}
\caption{Dynamic evolution of the Causal Effective Dimension ($q_t^*$). Shaded areas indicate historical stress episodes.}
\label{fig:qstar_dynamics}
\end{figure}

During liquidity crises, 80.8\% of observations exhibit $q^* \le 4$ with modal $q^* = 2$, compared to valuation repricing episodes where modal dimension is $q^* = 4$ (mean $q^* = 4.86$, with only 48.9\% exhibiting $q^* \le 4$). In baseline calm periods, the median dimension is 3 (mean $= 3.45$). Evaluating the correlation between $q^*$ and static Effective Rank across all 4,346 windows yields a Spearman rank correlation of $\rho_S = +0.1959$ (Pearson $\rho = +0.2262$), indicating that transition-based dimensionality is distinct from static covariance concentration.

\subsection{Linear Association with Conventional Systemic Risk Metrics}
To evaluate the relationship between $\mathrm{CEFI}_t$ and standard market metrics, we estimate:
\begin{equation}
\mathrm{CEFI}_t = \alpha + \beta_1 RV_t + \beta_2 \bar{\rho}_t + \beta_3 ER_t + \beta_4 DY_t + u_t
\end{equation}
where $RV_t$ is annualized realized volatility, $\bar{\rho}_t$ is average pairwise correlation, $ER_t = \exp(-\sum_i p_i \ln p_i)$ is effective rank \citep{roy2007effective}, and $DY_t$ is Diebold-Yilmaz total connectedness \citep{diebold2012better} (see Appendix A for formal definitions).

\begin{table}[htbp]
\centering
\caption{Linear Association Between $\mathrm{CEFI}_t$ and Conventional Systemic Risk Benchmarks (HAC Newey-West $L=40$)}
\label{tab:benchmarks_h4}
\resizebox{\textwidth}{!}{
\begin{tabular}{lccccc}
\toprule
\textbf{Benchmark Proxy} & \textbf{Pairwise Corr} & \textbf{Multivariate } $\mathbf{\beta}$ & \textbf{HAC SE} & $\mathbf{t\text{-stat}}$ & $\mathbf{p\text{-val}}$ \\
\midrule
Realized Volatility ($RV$)          & +0.628 & +1.0623 & (0.2359) & +4.50 & $6.67 \times 10^{-6}$ \\
Average Correlation ($\bar{\rho}$)  & +0.808 & +3.1668 & (0.4576) & +6.92 & $4.50 \times 10^{-12}$ \\
Effective Rank ($ER$)               & -0.761 & +0.0500 & (0.0187) & +2.68 & $7.40 \times 10^{-3}$ \\
Diebold-Yilmaz Index ($DY$)         & +0.740 & -0.0133 & (0.0059) & -2.27 & $2.30 \times 10^{-2}$ \\
\midrule
\multicolumn{6}{l}{$R^2 = 67.77\% \qquad \text{Unexplained Linear Variance} = 32.23\% \qquad N = 4,346$} \\
\bottomrule
\end{tabular}
}
\end{table}

Standard systemic risk proxies account for 67.77\% of the in-sample linear variation in $\mathrm{CEFI}_t$, while 32.23\% remains unexplained by this linear specification. Collinearity diagnostics yield $\text{VIF}(RV) = 3.20$, $\text{VIF}(\bar{\rho}) = 27.30$, $\text{VIF}(ER) = 13.77$, and $\text{VIF}(DY) = 15.11$ (condition number $= 11.80$). Residualized $\mathrm{CEFI}_t = \mathrm{CEFI}_t - \hat{\mathbb{E}}[\mathrm{CEFI}_t \mid RV, \bar{\rho}, ER, DY]$ retains positive means during GFC (+0.039) and COVID (+0.277) versus negative means during Dot-Com (-0.059), indicating that residualized $\mathrm{CEFI}_t$ retains crisis-associated episode-level variation after linear adjustment for conventional proxies.

\section{Robustness and Cross-Method Theoretical Benchmarking}
\label{sec:robustness}

\subsection{Cross-Method Benchmarking Against Published Literature}
Table \ref{tab:cross_framework} and Figure \ref{fig:cross_framework} report cross-method concordance across 870 historical rolling windows.

\begin{table}[htbp]
\centering
\caption{Cross-Method Benchmarking Against Published Literature (1992--2026, 870 Rolling Slices)}
\label{tab:cross_framework}
\resizebox{\textwidth}{!}{
\begin{tabular}{lcccc}
\toprule
\textbf{Theoretical Benchmark} & \textbf{Pearson } $\mathbf{\rho}$ & \textbf{Spearman } $\mathbf{\rho_S}$ & \textbf{Exact } $\mathbf{q^*}$ \textbf{Match} & $\mathbf{q^*}$ \textbf{Match (}$\pm 1$\textbf{)} \\
\midrule
Liu et al. (2024) Uniform $\Delta \mathcal{J}$ & 0.837 & 0.806 & 3.4\% & 5.1\% \\
Liu et al. (PRE 2025) SVD Emergence           & 0.832 & 0.764 & 75.4\% & \textbf{88.7\%} \\
\bottomrule
\end{tabular}
}
\end{table}

\begin{figure}[htbp]
\centering
\includegraphics[width=0.90\textwidth]{reports/figures/figure4_theoretical_benchmarking.pdf}
\caption{Cross-method benchmarking across 870 historical slices: (a) $\mathrm{CEFI}_t$ vs. exact $\Delta \mathcal{J}$ of \citet{liu2024exact}; (b) $\mathrm{CEFI}_t$ vs. analytical SVD emergence of \citet{liu2025singular}.}
\label{fig:cross_framework}
\end{figure}

The Stiefel manifold index is positively associated with the continuous formulation $\Delta \mathcal{J}$ of \citet{liu2024exact} (Pearson $\rho = 0.837$) and the analytical SVD emergence of \citet{liu2025singular} ($\rho = 0.832$). While \citet{liu2024exact} selects higher macro dimensions due to bounded uniform interventions, the SVD formulation yields 88.7\% dimensional agreement within $\pm 1$ dimension, providing cross-method evidence for recurrent low-dimensional organization.

\subsection{Cross-Universe Robustness (FF49)}
Replicating the pipeline across all $q \in \{1, \dots, 48\}$ on the 49 Industry Portfolios ($p=49$) yields a mean $\mathrm{CEFI}$ of 1.0585 and a modal dimension of $q^* = 3$, with 71.72\% of historical windows exhibiting $q^* \le 4$. In matched null inference during COVID 2020, $\mathrm{CEFI}_{\text{obs}} = 1.6350 \ge Q_{95}(H_0^{\text{static}}) = 1.6243$. In finite-sample surrogate testing with $B=100$, exactly 5 surrogate realizations equal or exceed the observed value ($k=5$), yielding an empirical $p$-value of $p_{\text{emp}} = \frac{1+5}{100+1} = 0.0594$; the finite-sample test therefore does not reject at the 5\% significance level, serving as an exploratory cross-universe check.

\subsection{Sensitivity to Scale Parameter \texorpdfstring{$\kappa$}{kappa} and Window Size \texorpdfstring{$W$}{W}}
Varying the intervention strength across $\kappa \in \{0.25, 0.5, 1.0, 2.0, 4.0\}$ preserves rank correlations ($\rho_S \ge 0.718$), low modal dimensions ($q^* \le 3$), and crisis significance ($\beta_{\text{Liq}} > 0, t \ge 2.92$). Increasing estimation windows to $W = 750$ and $W = 1000$ days yields rank correlations of $\rho_S = 0.726$ and $\rho_S = 0.593$, with peak cross-correlation occurring at a phase shift of $\ell = -90$ days ($\rho = 0.683$), reflecting mechanical rolling window smoothing.

\section{Discussion}
\label{sec:discussion}

\subsection{Economic Interpretation of Financial Causal Emergence}
A positive $\mathrm{CEFI}_t$ indicates that an optimal macroscopic coarse-graining of the financial system possesses a higher causal transition density than the disaggregated micro-level assets. In economic terms, when individual industries experience elevated idiosyncratic noise, a properly oriented lower-dimensional projection can eliminate mutually cancelling noise components while concentrating collective intertemporal momentum.

\subsection{Disentangling Contemporaneous Co-Movement from Lagged Dynamical Structure}
Our four-tier null hierarchy clarifies the distinction between static correlation and dynamic causal organization. In the 2005 calm-market benchmark, observed $\mathrm{CEFI}_t$ is consistent with contemporaneous correlation and own-lag persistence ($p_{\text{emp}} = 0.6215, p_{\text{Holm}} = 0.6215$). In systemic liquidity crises (e.g., the 2008 GFC peak), $\mathrm{CEFI}_t$ significantly exceeds both the static correlation benchmark ($p = 0.0014$) and cross-lag isolation ($p = 0.0001$), indicating the presence of structured intertemporal organization that cannot be explained by static covariance alone.

\subsection{Inter-Crisis Episodes (2011--2013 Debt Ceiling and Sovereign Debt)}
As shown in Figure \ref{fig:cefi_dynamics}, $\mathrm{CEFI}_t$ also displays secondary peaks during the 2011 U.S. sovereign debt downgrade and the European sovereign debt crisis. These peaks coincide with known macro-financial stress events. However, because matched-null inference was not conducted for every historical window, their interpretation as dynamic cross-sectoral organization remains descriptive.

\subsection{Complementarity with Macroprudential Indicators}
These results suggest that causal emergence may provide a complementary information-theoretic lens for monitoring changes in systemic market organization, capturing shifts in collective transition structure that supplement traditional static correlation and volatility indices.

\section{Limitations}
\label{sec:limitations}

Several methodological limitations should be noted:
\begin{enumerate}
    \item \textbf{Optimizer Search Budget Sensitivity}: While the Stiefel optimization exhibits substantial index concordance (Pearson $\rho = 0.8913$; Spearman $\rho_S = 0.8062$) and high $\pm 1$ dimensional consistency (84.0\%) relative to a high-budget reference configuration, exact objective values and dimension selections remain partially sensitive to the finite multistart optimization budget.
    \item \textbf{Linear Gaussian Transition Dynamics}: Micro-dynamics are modeled through a linear VAR(1) with Gaussian innovations. Non-linear spillover dynamics, volatility clustering, and heavy-tailed co-jumps are not explicitly parameterized in the transition operator.
    \item \textbf{Reduced-Form Representation}: The transition matrix $\mathbf{A}$ is estimated in reduced form. Interventions represent counterfactual Gaussian perturbations under the fitted dynamical model and should not be interpreted as structural policy shocks.
    \item \textbf{Generated-Statistic Uncertainty}: $\mathrm{CEFI}_t$ is a generated statistic obtained from first-stage parameter estimates ($\hat{\mathbf{A}}, \boldsymbol{\Sigma}_\varepsilon, \mathbf{W}, q^*$). Downstream HAC regressions do not fully propagate first-stage estimation uncertainty.
    \item \textbf{Aggregation Level}: The empirical analysis relies on industry portfolios rather than individual financial institutions, reflecting cross-sectoral market organization rather than interbank counterparty networks.
    \item \textbf{Sample Size of Historical Crises}: The number of major systemic dislocations in modern market history is inherently limited. Event-study inference should be interpreted as within-sample regime evidence across historical episodes rather than asymptotic identification of crisis taxonomies.
    \item \textbf{Intervention Family Dependence}: Although empirical conclusions remain stable across the evaluated $\kappa$ range, $\mathrm{CEFI}_t$ remains conditional on the chosen family of isotropic Gaussian interventions.
    \item \textbf{Absence of Early-Warning Validation}: The empirical design evaluates concurrent market organization during stress regimes; no early-warning or predictive lead-time capability is claimed or validated.
\end{enumerate}

\section{Conclusion}
\label{sec:conclusion}

This paper has developed and evaluated continuous-state Causal Emergence as an empirical measure of dynamic cross-sectoral market organization during systemic stress. By formulating scale-invariant Effective Information on Stiefel manifolds and benchmarking against a hierarchy of matched surrogate null models ($B=9,999$ for primary nulls and $B=999$ for auxiliary nulls), we find that causal emergence varies across market states and historical episodes. While the 2005 calm-market benchmark is consistent with univariate persistence and contemporaneous covariance, systemic liquidity and contagion crises exhibit elevated macroscopic organization and a concentration of causal effective dimensionality toward low-dimensional macro representations (combined modal $q^* = 2$, with 80.8\% of observations having $q^* \le 4$, compared with modal $q^* = 4$ and 48.9\% having $q^* \le 4$ during valuation repricing episodes). These results indicate that continuous causal emergence provides an informative framework for monitoring collective market organization during financial distress.

\appendix

\section{Benchmark Construction and Definitions}
\label{app:benchmarks}
Conventional systemic risk proxies are constructed over each rolling window $X_{t-W+1:t}$ as follows:
\begin{enumerate}
    \item \textbf{Realized Volatility ($RV_t$)}: Annualized standard deviation of the equal-weighted portfolio return series, $RV_t = \sqrt{252} \cdot \operatorname{std}(\bar{x}_t)$.
    \item \textbf{Average Correlation ($\bar{\rho}_t$)}: Mean off-diagonal element of the $p \times p$ sample correlation matrix $\mathbf{R}_t$: $\bar{\rho}_t = \frac{2}{p(p-1)} \sum_{i < j} R_{ij,t}$.
    \item \textbf{Effective Rank ($ER_t$)}: Spectral entropy exponentiation of normalized covariance eigenvalues \citep{roy2007effective}:
    \begin{equation}
    ER_t = \exp\left( -\sum_{i=1}^p p_i \ln p_i \right), \qquad p_i = \frac{\lambda_i(\boldsymbol{\Sigma}_{x,t})}{\sum_{j=1}^p \lambda_j(\boldsymbol{\Sigma}_{x,t})}
    \end{equation}
    \item \textbf{Diebold-Yilmaz Connectedness Index ($DY_t$)}: Total volatility spillover index derived from generalized forecast error variance decompositions of a VAR(1) with forecast horizon $H = 10$ days \citep{diebold2012better}.
\end{enumerate}

\section*{Declaration of Competing Interests}
The author declares that they have no known competing financial interests or personal relationships that could have appeared to influence the work reported in this paper.

\section*{Data and Code Availability}
All datasets used in this study (Kenneth French Industry Portfolios) are publicly available. Replication code and estimation routines will be deposited in an open-access repository upon publication.

\section*{Declaration of Generative AI and AI-Assisted Technologies in the Writing Process}
During the preparation of this work, the author used generative AI tools (OpenAI ChatGPT and Google Gemini) to assist with code optimization, numerical verification scripting, LaTeX typesetting, and language refinement. All research conceptualization, mathematical formulations, econometric methodology, empirical results, and interpretations were directed, verified, and validated by the author, who assumes full scientific responsibility for the content of this manuscript.

\bibliographystyle{elsarticle-harv}
\begin{thebibliography}{28}
\bibitem[Absil et~al., 2008]{absil2008optimization}
Absil, P.-A., Mahony, R., Sepulchre, R., 2008.
\newblock Optimization Algorithms on Matrix Manifolds.
\newblock Princeton University Press, Princeton, NJ. \url{https://doi.org/10.1515/9781400830244}

\bibitem[Acharya et~al., 2017]{acharya2017measuring}
Acharya, V. V., Pedersen, L. H., Philippon, T., Richardson, M., 2017.
\newblock Measuring Systemic Risk.
\newblock \emph{Review of Financial Studies} 30(1), 2--47. \url{https://doi.org/10.1093/rfs/hhw088}

\bibitem[Adrian et~al., 2016]{adrian2016covar}
Adrian, T., Brunnermeier, M. K., 2016.
\newblock CoVaR.
\newblock \emph{American Economic Review} 106(7), 1705--1741. \url{https://doi.org/10.1257/aer.20120555}

\bibitem[Ahelegbey et~al., 2022]{ahelegbey2022network}
Ahelegbey, D. F., Cerchiello, P., Scaramozzino, R., 2022.
\newblock Network based evidence of the financial impact of Covid-19 pandemic.
\newblock \emph{International Review of Financial Analysis} 81, 102101. \url{https://doi.org/10.1016/j.irfa.2022.102101}

\bibitem[Bekiros et~al., 2017]{bekiros2017information}
Bekiros, S., Nguyen, D. K., Uddin, G. S., Sjo, B., 2017.
\newblock Information diffusion, cluster formation and entropy-based network dynamics in equity and commodity markets.
\newblock \emph{European Journal of Operational Research} 256(3), 945--961. \url{https://doi.org/10.1016/j.ejor.2016.06.052}

\bibitem[Billio et~al., 2012]{billio2012econometric}
Billio, M., Getmansky, M., Lo, A. W., Pelizzon, L., 2012.
\newblock Econometric measures of connectedness and systemic risk in the finance and insurance sectors.
\newblock \emph{Journal of Financial Economics} 104(3), 535--559. \url{https://doi.org/10.1016/j.jfineco.2011.12.010}

\bibitem[Bouri et~al., 2021]{bouri2021return}
Bouri, E., Cepni, O., Gabauer, D., Gupta, R., 2021.
\newblock Return connectedness across asset classes around the COVID-19 outbreak.
\newblock \emph{International Review of Financial Analysis} 73, 101646. \url{https://doi.org/10.1016/j.irfa.2020.101646}

\bibitem[Brownlees et~al., 2017]{brownlees2017srisk}
Brownlees, C., Engle, R. F., 2017.
\newblock SRISK: A conditional capital shortfall measure of systemic risk.
\newblock \emph{Review of Financial Studies} 30(1), 48--79. \url{https://doi.org/10.1093/rfs/hhw060}

\bibitem[Diebold et~al., 2012]{diebold2012better}
Diebold, F. X., Yilmaz, K., 2012.
\newblock Better to give than to receive: Predictive directional measurement of volatility spillovers.
\newblock \emph{International Journal of Forecasting} 28(1), 57--66. \url{https://doi.org/10.1016/j.ijforecast.2011.02.006}

\bibitem[Diebold et~al., 2014]{diebold2014network}
Diebold, F. X., Yilmaz, K., 2014.
\newblock On the network topology of variance decompositions: Measuring the connectedness of financial firms.
\newblock \emph{Journal of Econometrics} 182(1), 119--134. \url{https://doi.org/10.1016/j.jeconom.2014.04.012}

\bibitem[Dimpfl et~al., 2013]{dimpfl2013using}
Dimpfl, T., Peter, F. J., 2013.
\newblock Using transfer entropy to measure dynamic information flows between markets.
\newblock \emph{Journal of International Financial Markets, Institutions and Money} 27, 273--294. \url{https://doi.org/10.1016/j.intfin.2013.09.003}

\bibitem[Edelman et~al., 1998]{edelman1998geometry}
Edelman, A., Arias, T. A., Smith, S. T., 1998.
\newblock The geometry of algorithms with orthogonality constraints.
\newblock \emph{SIAM Journal on Matrix Analysis and Applications} 20(2), 303--353. \url{https://doi.org/10.1137/S0895479895290954}

\bibitem[Hoel, 2017]{hoel2017when}
Hoel, E. P., 2017.
\newblock When the map is better than the territory.
\newblock \emph{Entropy} 19(5), 188. \url{https://doi.org/10.3390/e19050188}

\bibitem[Hoel et~al., 2013]{hoel2013quantifying}
Hoel, E. P., Albantakis, L., Tononi, G., 2013.
\newblock Quantifying causal emergence shows that macro can beat micro.
\newblock \emph{Proceedings of the National Academy of Sciences} 110(49), 19790--19795. \url{https://doi.org/10.1073/pnas.1314922110}

\bibitem[Klein et~al., 2020]{klein2020emergence}
Klein, B., Hoel, E., 2020.
\newblock The emergence of informative higher scales in complex networks.
\newblock \emph{Complexity} 2020, 8932526. \url{https://doi.org/10.1155/2020/8932526}

\bibitem[Kritzman et~al., 2011]{kritzman2011principal}
Kritzman, M., Li, Y., Page, S., Rigobon, R., 2011.
\newblock Principal components as a measure of systemic risk.
\newblock \emph{Journal of Portfolio Management} 37(4), 112--126. \url{https://doi.org/10.3905/jpm.2011.37.4.112}

\bibitem[Ledoit et~al., 2004]{ledoit2004well}
Ledoit, O., Wolf, M., 2004.
\newblock A well-conditioned estimator for large-dimensional covariance matrices.
\newblock \emph{Journal of Multivariate Analysis} 88(2), 365--411. \url{https://doi.org/10.1016/S0047-259X(03)00096-4}

\bibitem[Liu et~al., 2024]{liu2024exact}
Liu, K., Yuan, B., Zhang, J., 2024.
\newblock An exact theory of causal emergence for linear stochastic iteration systems.
\newblock \emph{Entropy} 26(8), 618. \url{https://doi.org/10.3390/e26080618}

\bibitem[Liu et~al., 2025]{liu2025singular}
Liu, K., Pan, L., Wang, Z., Yang, M., Yuan, B., Zhang, J., 2025.
\newblock Singular-value-decomposition-based causal emergence for Gaussian iterative systems.
\newblock \emph{Physical Review E} 112(5), 054225. \url{https://doi.org/10.1103/PhysRevE.112.054225}

\bibitem[Mensi et~al., 2021]{mensi2021spillovers}
Mensi, W., Hernandez, J. A., Yoon, S.-M., Vo, X. V., Kang, S. H., 2021.
\newblock Spillovers and connectedness between major precious metals and major currency markets: The role of frequency factor.
\newblock \emph{International Review of Financial Analysis} 74, 101672. \url{https://doi.org/10.1016/j.irfa.2021.101672}

\bibitem[Newey et~al., 1987]{newey1987hypothesis}
Newey, W. K., West, K. D., 1987.
\newblock A simple, positive semi-definite, heteroskedasticity and autocorrelation consistent covariance matrix.
\newblock \emph{Econometrica} 55(3), 703--708. \url{https://doi.org/10.2307/1913610}

\bibitem[Politis et~al., 1994]{politis1994stationary}
Politis, D. N., Romano, J. P., 1994.
\newblock The stationary bootstrap.
\newblock \emph{Journal of the American Statistical Association} 89(428), 1303--1313. \url{https://doi.org/10.1080/01621459.1994.10476870}

\bibitem[Rosas et~al., 2020]{rosas2020reconciling}
Rosas, F. E., Mediano, P. A., Jensen, H. J., Seth, A. K., Barrett, A. B., Carhart-Harris, R. L., Bor, D., 2020.
\newblock Reconciling emergences: An information-theoretic approach to identify causal emergence in multivariate data.
\newblock \emph{PLoS Computational Biology} 16(12), e1008579. \url{https://doi.org/10.1371/journal.pcbi.1008579}

\bibitem[Roy et~al., 2007]{roy2007effective}
Roy, O., Vetterli, M., 2007.
\newblock The effective rank: A measure of effective dimensionality.
\newblock 15th European Signal Processing Conference (EUSIPCO 2007), 606--610.

\bibitem[Schreiber, 2000]{schreiber2000measuring}
Schreiber, T., 2000.
\newblock Measuring information transfer.
\newblock \emph{Physical Review Letters} 85(2), 461--464. \url{https://doi.org/10.1103/PhysRevLett.85.461}

\bibitem[Yang et~al., 2025]{yang2025finding}
Yang, M., Wang, Z., Liu, K., Rong, Y., Yuan, B., Zhang, J., 2025.
\newblock Finding emergence in data by maximizing effective information.
\newblock \emph{National Science Review} 12(1), nwae279. \url{https://doi.org/10.1093/nsr/nwae279}

\bibitem[Yang et~al., 2023]{yang2023modeling}
Yang, L., Hamori, S., 2023.
\newblock Modeling the global sovereign credit network under climate change.
\newblock \emph{International Review of Financial Analysis} 87, 102618. \url{https://doi.org/10.1016/j.irfa.2023.102618}

\bibitem[Yousaf et~al., 2022]{yousaf2022linkages}
Yousaf, I., Nekhili, R., Gubareva, M., 2022.
\newblock Linkages between DeFi assets and conventional currencies: Evidence from the COVID-19 pandemic.
\newblock \emph{International Review of Financial Analysis} 81, 102082. \url{https://doi.org/10.1016/j.irfa.2022.102082}
\end{thebibliography}

\end{document}
